\documentclass[a4paper,11pt]{article}
\usepackage[T1]{fontenc}
\usepackage[ngerman]{babel}
\usepackage[margin=2cm]{geometry}
\usepackage{amsmath}
\usepackage{booktabs}
\usepackage{siunitx}
\usepackage{threeparttable}
\sisetup{per-mode=symbol}

\DeclareMathOperator{\sinc}{sinc}
\DeclareMathOperator{\sign}{sign}

\begin{document}

\begin{table}[htbp]
    \centering
    \small
    \renewcommand{\arraystretch}{1.6}
    \setlength{\tabcolsep}{4pt}
    \begin{threeparttable}
        \caption{Einheiten spezieller Funktionen im Zeit- und Frequenzbereich}
        \label{tab:einheiten}
        \begin{tabular}{@{}l l l l l@{}}
            \toprule
            \textbf{Funktion} & \textbf{Zeitbereich}                                                       & \textbf{Einheit}                           & \textbf{Frequenzbereich} & \textbf{Einheit} \\
            \midrule
            \multicolumn{5}{@{}l}{\textit{Dimensionslose Grundfunktionen}}                                                                                                                            \\
            Dirac-Impuls
                              & $\delta(t)$                                                                & $\si{\per\second}$
                              & $1$                                                                        & $1$                                                                                      \\
            Konstante
                              & $1$                                                                        & $1$
                              & $\delta(f)$                                                                & $\si{\per\hertz} = \si{\second}$                                                         \\
            Komplexe Schwingung
                              & $e^{j2\pi f_0 t}$                                                          & $1$
                              & $\delta(f-f_0)$                                                            & $\si{\second}$                                                                           \\
            Einheitssprung
                              & $u(t)$                                                                     & $1$
                              & $\frac{1}{2}\delta(f) + \frac{1}{j2\pi f}$                                 & $\si{\second}$                                                                           \\
            Signumfunktion
                              & $\sign(t)$                                                                 & $1$
                              & $\frac{1}{j\pi f}$                                                         & $\si{\second}$                                                                           \\
            Rechteckimpuls
                              & $u(t+\frac{T}{2}) - u(t-\frac{T}{2})$                                      & $1$
                              & $T\sinc(\pi f T)$                                                          & $\si{\second}$                                                                           \\
            Dirac-Impulsfolge
                              & $\sum_{n} \delta(t-nT)$                                                    & $\si{\per\second}$
                              & $\frac{1}{T}\sum_{k} \delta\!\left(f-\frac{k}{T}\right)$                   & $1$                                                                                      \\
            \midrule
            \multicolumn{5}{@{}l}{\textit{Signale (Beispiel: Spannung)}}                                                                                                                              \\
            Energiesignal
                              & $x(t)$                                                                     & $\si{\volt}$
                              & $X(f)$                                                                     & $\si{\volt\per\hertz} = \si{\volt\second}$                                               \\
            Periodisches Signal\tnote{a}
                              & $\sum_{k} c_k\, e^{j2\pi k f_0 t}$                                         & $\si{\volt}$
                              & $\sum_{k} c_k\, \delta(f-kf_0)$                                            & $\si{\volt\per\hertz}$                                                                   \\
            Energie-AKF
                              & $\varphi_{xx}(\tau) = \int x(t)\,x(t+\tau)\,\mathrm{d}t$                   & $\si{\volt\squared\second}$
                              & $|X(f)|^2$                                                                 & $\si{\volt\squared\second\per\hertz}$                                                    \\
            \midrule
            \multicolumn{5}{@{}l}{\textit{LTI-System (Ein- und Ausgang in \si{\volt})}\tnote{b}}                                                                                                      \\
            Impulsantwort
                              & $h(t)$                                                                     & $\si{\per\second}$
                              & $H(f)$                                                                     & $1$                                                                                      \\
            Sprungantwort
                              & $s(t) = \int_{-\infty}^{t} h(\tau)\,\mathrm{d}\tau$                        & $1$
                              & $\frac{H(f)}{j2\pi f} + \frac{1}{2}H(0)\,\delta(f)$                        & $\si{\second}$                                                                           \\
            Ausgangssignal
                              & $y(t) = x(t) * h(t)$                                                       & $\si{\volt}$
                              & $Y(f) = X(f)\,H(f)$                                                        & $\si{\volt\per\hertz}$                                                                   \\
            \midrule
            \multicolumn{5}{@{}l}{\textit{Zeitdiskret (DTFT, $\Omega$ dimensionslos)}\tnote{c}}                                                                                                       \\
            Einheitsimpuls
                              & $\delta[n]$                                                                & $1$
                              & $1$                                                                        & $1$                                                                                      \\
            Sprungfolge
                              & $u[n]$                                                                     & $1$
                              & $\frac{1}{1-e^{-j\Omega}} + \pi\sum_{k}\delta(\Omega-2\pi k)$              & $1$                                                                                      \\
            Signal
                              & $x[n]$                                                                     & $\si{\volt}$
                              & $X(\Omega) = \sum_{n} x[n]\,e^{-j\Omega n}$                                & $\si{\volt}$                                                                             \\
            Impulsantwort
                              & $h[n]$                                                                     & $1$
                              & $H(\Omega)$                                                                & $1$                                                                                      \\
            Sprungantwort
                              & $g[n] = \sum_{m=-\infty}^{n} h[m]$                                         & $1$
                              & $\frac{H(\Omega)}{1-e^{-j\Omega}} + \pi H(0)\sum_{k}\delta(\Omega-2\pi k)$ & $1$                                                                                      \\
            \bottomrule
        \end{tabular}
        \begin{tablenotes}\footnotesize
            \item[] Faustregel: $[X(f)] = [x(t)]\cdot\si{\second}$, da das Fourier-Integral über $\mathrm{d}t$ läuft.
            In $\omega$-Darstellung bleiben die SI-Einheiten gleich (rad ist dimensionslos), es gilt aber $\delta(f) = 2\pi\,\delta(\omega)$.
            \item[a] Die Fourier-Koeffizienten $c_k$ und damit die Impulsgewichte im Spektrum haben die Einheit \si{\volt}.
            \item[b] Allgemein (Eingang in $[x]$, Ausgang in $[y]$):
            $[h(t)] = \frac{[y]}{[x]\cdot\si{\second}}$, \quad $[H(f)] = [s(t)] = \frac{[y]}{[x]}$, \quad $[S(f)] = \frac{[y]}{[x]}\cdot\si{\second}$.
            \item[c] Die DTFT ist eine Summe ohne $\mathrm{d}t$, daher keine zusätzliche Einheit.
            Wird $h[n]$ aus einer zeitkontinuierlichen Impulsantwort gewonnen, gilt $h[n] = T_\mathrm{s}\, h(nT_\mathrm{s})$ (dimensionslos).
        \end{tablenotes}
    \end{threeparttable}
\end{table}

\end{document}